\subsection{Tensor products, spherical harmonics, and the Wigner-Eckart viewpoint}
\label{subsec:tensor_products_background}

The basic algebraic operation of an equivariant graph neural network is the
tensor product of irreps. If two feature channels transform according to
$D^{(\ell_1,p_1)}$ and $D^{(\ell_2,p_2)}$, then their tensor product
decomposes as
\begin{equation}
D^{(\ell_1,p_1)} \otimes D^{(\ell_2,p_2)}
=
\bigoplus_{L=|\ell_1-\ell_2|}^{\ell_1+\ell_2}
D^{(L,p_1 p_2)}.
\end{equation}
This Clebsch-Gordan structure determines which output irreps may be produced by
combining two input channels. In practical terms, it is the algebraic rule that
governs message passing, nonlinear mixing of hidden features, and the final
projection from latent features to matrix coefficients.

Edge geometry is injected through spherical harmonics
$Y_{\ell m}(\hat{\bm{r}}_{ij})$, evaluated on normalized relative displacement
directions. Because spherical harmonics transform irreducibly under rotations,
they provide the canonical angular basis for equivariant edge messages. A
typical equivariant interaction therefore combines learned node or edge
features with spherical-harmonic features through tensor products, while radial
functions modulate the interaction strength as a function of distance.

The same representation-theoretic structure appears in matrix elements of
operators between angular-momentum-carrying orbitals. For a spherical tensor
operator $T^{(k)}_q$, the Wigner-Eckart theorem states that
\begin{equation}
\langle \alpha \ell m | T^{(k)}_q | \alpha' \ell' m' \rangle
=
\langle \ell' m'; k q | \ell m \rangle
\langle \alpha \ell \| T^{(k)} \| \alpha' \ell' \rangle,
\end{equation}
that is, the dependence on magnetic quantum numbers is separated from the
reduced matrix element. Mandala uses this viewpoint operationally: the angular
coupling structure is fixed analytically by the irrep mapping, while the neural
network learns the reduced coefficients associated with each orbital pair and
chemical environment. This is precisely why the matrix-block prediction problem
can be posed naturally in an irrep basis rather than directly in Cartesian
matrix entries.
